\section{Discussion and Limitations}
\label{sec:discussion}

The results presented in Section~\ref{sec:evaluation} demonstrate that VulnRAG significantly enhances the intelligence gathering capabilities of autonomous security agents. However, the deployment of such systems in enterprise environments raises several critical considerations regarding resilience, ethics, and technical constraints.

\subsection{Operational Resilience and Reliability}
A primary advantage of VulnRAG’s modular agentic architecture is its resilience to provider-specific failures. By utilizing a provider-agnostic self-correction loop, the system can adapt to degraded performance or outages in individual LLM backends. Furthermore, our complexity-aware routing ensures that the system remains economically viable during large-scale assessments, where the cost of flagship model inference would otherwise be prohibitive. However, the reliance on external APIs (both for LLM inference and NVD synchronization) introduces a degree of "cascading latency." While our tiered routing mitigates this, real-time assessment of ultra-low latency environments remains a challenge for future optimization.

\subsection{Ethical Considerations and Dual-Use Mitigation}
The development of highly autonomous vulnerability discovery systems is inherently subject to the "dual-use" dilemma. While VulnRAG is designed as a defensive tool for security auditors and red-teamers, its capabilities could theoretically be misappropriated for malicious reconnaissance. To mitigate this risk, we emphasize the following safety guardrails:
\begin{itemize}
    \item \textbf{Restricted Execution}: The VulnRAG framework should be deployed within controlled environments with strict egress filtering.
    \item \textbf{Human-in-the-Loop (HITL)}: While retrieval and initial reasoning are autonomous, we recommend a HITL approach for the final validation of exploit-ability to prevent unintentional damage to production systems.
    \item \textbf{Transparency}: VulnRAG maintains detailed "Reasoning Traces," allowing human auditors to inspect the logic behind every retrieved intelligence result and attack path.
\end{itemize}

\subsection{Failure Mode Analysis and Hallucinations}
Despite the robustness of the Evaluator-Reformulator loop, the system remains susceptible to two primary failure modes:
\begin{itemize}
    \item \textbf{Silent Routing Errors}: A low-complexity (Flash) model may return a "confident" but factually incorrect response, preventing escalation to the Reasoning tier.
    \item \textbf{Graph Over-Exploration}: In rare cases, the Graph Explorer may prioritize "noisy" edges in legacy software, leading to context saturation. 
\end{itemize}
Future work will explore the use of \textbf{Logit-Bias} constraints to bound the creativity of the Reformulator and ensure strict adherence to structured vulnerability schemas.

\subsection{Limitations and Technical Constraints}
Despite its innovations, VulnRAG has several limitations:
\begin{itemize}
    \item \textbf{NVD Rate Limiting}: The \texttt{nvd\_sync\_service.py} is constrained by the NIST NVD API 2.0 rate limits. During high-intensity vulnerability "storms" (e.g., major library zero-days), the system may face synchronization delays.
    \item \textbf{Graph Depth Complexity}: Structural discovery via \texttt{cve\_graph.py} is currently capped at a depth of 3 hops to prevent "graph explosion" and excessive token consumption during synthesis. This may miss extremely deep, multi-link transitive dependencies.
    \item \textbf{Semantic Ambiguity}: While agentic reformulation handles most cases, extreme ambiguity in product naming (e.g., generic terms like "OS" or "Server") still results in lower precision that requires manual oversight.
\end{itemize}

Addressing these limitations, particularly through the integration of local, high-performance Knowledge Graphs (KGs) and specialized small-parameter security models, represents the next frontier for VulnRAG’s evolution.
